\documentclass[DIN, pagenumber=false, fontsize=11pt, parskip=half]{scrartcl}
\usepackage[UTF8]{ctex}  % 中文支持
% \usepackage{ngerman}
\usepackage[utf8]{inputenc}
\usepackage[T1]{fontenc}
\usepackage{textcomp}
\usepackage{xcolor} % Required for custom red

\usepackage{geometry} % Custom margins

\usepackage{tocloft} % Modifying the table of contents

% for matlab code
% bw = blackwhite - optimized for print, otherwise source is colored
\usepackage[framed,numbered,bw]{mcode}
\usepackage[most]{tcolorbox}
% for other code
%\usepackage{listings}

\setlength{\parindent}{0em}

% set section in CM
\setkomafont{section}{\normalfont\bfseries\Large}

\newcommand{\mytitle}[1]{{\noindent\Large\textbf{#1}}}
\newcommand{\mysection}[1]{\textbf{\section*{#1}}}
\newcommand{\mysubsection}[2]{\romannumeral #1) #2}

\newtcolorbox{parchmentquote}{
    enhanced,
    colback=orange!5!white,   % 浅米黄色背景
    colframe=orange!50!black,
    boxrule=0.5pt,
    arc=2pt,                  % 细微圆角
    boxsep=5pt,
    before skip=10pt,
    after skip=10pt,
    fontupper=\itshape,
    borderline west={2pt}{0pt}{orange!80!black} % 加厚的左装饰线
}
%===================================
\begin{document}

\noindent\textbf{课下练习} \hfill \textbf{online course}\\
primary school \hfill Liu\\

\mytitle{分数约分与通分 \hfill \today}
\newline
\begin{parchmentquote}
	应用分数基本性质，可以实现分数的约分与通分
\end{parchmentquote}


%===================================
\section{简单约分与通分运算}

\subsection{化成最简分数} 

\begin{enumerate}
    \item \(\frac{12}{18} = \underline{\qquad}\)
    \item \(\frac{45}{60} = \underline{\qquad}\)
    \item \(\frac{28}{42} = \underline{\qquad}\)
    \item \(\frac{36}{48} = \underline{\qquad}\)
    \item \(\frac{16}{40} = \underline{\qquad}\)
\end{enumerate}
\subsection{分数通分} 

\noindent \textbf{练习：在括号内填入合适的数字完成通分。}

\begin{enumerate}
    \item $\dfrac{2}{3} = \dfrac{2 \times (\quad)}{3 \times 4} = \dfrac{(\quad)}{12}$
    \item $\dfrac{17}{19}  = \dfrac{(\quad)}{152}$
    \item $\dfrac{13}{23}  = \dfrac{39}{(\quad)}$
    \item $\dfrac{5}{8}$ 和 $\dfrac{7}{12}$ 转成分母是24的分数：
    \[
    \frac{5}{8} = \frac{15}{24}, \quad \frac{7}{12} = \frac{(\quad)}{24}
    \]
     \item $\dfrac{5}{6}$ 和 $\dfrac{7}{9}$ 最小公倍数通分：
    \[
    \frac{5}{6} = \frac{(\quad)}{(\quad)}, \quad \frac{7}{9} = \frac{(\quad)}{(\quad)}
    \]
\end{enumerate}

\section{应用题}
\subsection{比较大小}
小明、小红和小华进行打字比赛。在相同的时间内，小明完成了打字任务的 $\frac{5}{8}$，小红完成了 $\frac{7}{12}$，小华完成了 $\frac{3}{4}$。请问谁的打字速度最快？
\subsection{信息提取}
一个分数的分子是最小的质数，分母比分子大6，求这个分数并化成最简形式。
\end{document}